%%%%%%%%%%%%%%%%%%%%%%%%%%%%%%%%%%%%%%%%%
% Academic Title Page
% LaTeX Template
% Version 2.0 (17/7/17)
%
% This template was downloaded from:
% http://www.LaTeXTemplates.com
%
% Original author:
% WikiBooks (LaTeX - Title Creation) with modifications by:
% Vel (vel@latextemplates.com)
%
% License:
% CC BY-NC-SA 3.0 (http://creativecommons.org/licenses/by-nc-sa/3.0/)
% 
% Instructions for using this template:
% This title page is capable of being compiled as is. This is not useful for 
% including it in another document. To do this, you have two options: 
%
% 1) Copy/paste everything between \begin{document} and \end{document} 
% starting at \begin{titlepage} and paste this into another LaTeX file where you 
% want your title page.
% OR
% 2) Remove everything outside the \begin{titlepage} and \end{titlepage}, rename
% this file and move it to the same directory as the LaTeX file you wish to add it to. 
% Then add \input{./<new filename>.tex} to your LaTeX file where you want your
% title page.
%
%%%%%%%%%%%%%%%%%%%%%%%%%%%%%%%%%%%%%%%%%

%----------------------------------------------------------------------------------------
%	PACKAGES AND OTHER DOCUMENT CONFIGURATIONS
%----------------------------------------------------------------------------------------

\documentclass[11pt,letterpaper ,twoside ,openright]{book}

\usepackage{mathpazo} % Palatino font
\usepackage{graphicx}

\newcommand\abstractname{Abstract}
\newenvironment{abstract}
{\cleardoublepage\null \vfill\begin{center}
\bfseries \abstractname \end{center}}
	{\vfill\null}

\usepackage{fancyhdr}
\pagestyle{empty}


\begin{document}

%----------------------------------------------------------------------------------------
%	TITLE PAGE
%----------------------------------------------------------------------------------------

\begin{titlepage} % Suppresses displaying the page number on the title page and the subsequent page counts as page 1
	\newcommand{\HRule}{\rule{\linewidth}{0.5mm}} % Defines a new command for horizontal lines, change thickness here
	
	\center
	
	%------------------------------------------------
	%	Headings
	%------------------------------------------------
	
	\textsc{\LARGE Royal Holloway, University of London}\\[1.5cm] % Main heading such as the name of your university/college
	
	\textsc{\Large Full Unit Project}\\[0.5cm] % Major heading such as course name
	
	\textsc{\large Final Project Report}\\[0.5cm] % Minor heading such as course title
	
	%------------------------------------------------
	%	Title
	%------------------------------------------------
	
	\HRule\\[0.4cm]
	
	{\huge\bfseries Cooperative Strategies in Multi-Agent Systems}\\[0.4cm] % Title of your document
	
	\HRule\\[1.5cm]
	
	%------------------------------------------------
	%	Author(s)
	%------------------------------------------------
	
	\begin{minipage}{0.4\textwidth}
		\begin{flushleft}
			\large
			\textit{Author}\\
			James \textsc{King} % Your name
		\end{flushleft}
	\end{minipage}
	~
	\begin{minipage}{0.4\textwidth}
		\begin{flushright}
			\large
			\textit{Supervisor}\\
			Prof. Kostas \textsc{Stathis} % Supervisor's name
		\end{flushright}
	\end{minipage}
	
	% If you don't want a supervisor, uncomment the two lines below and comment the code above
	%{\large\textit{Author}}\\
	%John \textsc{Smith} % Your name
	
	%------------------------------------------------
	%	Date
	%------------------------------------------------
	
	\vfill\vfill\vfill % Position the date 3/4 down the remaining page
	
	{\large\today} % Date, change the \today to a set date if you want to be precise
	
	%------------------------------------------------
	%	Logo
	%------------------------------------------------
	
	\vfill\vfill
	\includegraphics[width=0.2\textwidth]{Rhulcoa.png}\\[1cm] % Include a department/university logo - this will require the graphicx package
	 
	%----------------------------------------------------------------------------------------
	
	\vfill % Push the date up 1/4 of the remaining page
		
		
\end{titlepage}
%----------------------------------------------------------------------------------------

\section*{Acknowledgements}
\addtocounter{section}{1}
%I would like to thank my awesome girlfriend Kira cos she is amazing and the bestest ever.\\
%Plus she said to split up strategies into basic and advanced :P\\
Miguel Grinberg, Vince Knight, Nikoleta Glynetski, Owen Campbell, Annie Ogborne, Martin A. Nowak, Robert Axelrod

\tableofcontents
\listoffigures
\listoftables

\begin{abstract}
Hello this is my abstract\\

https://users.ece.cmu.edu/~koopman/essays/abstract.html
\end{abstract}

\chapter{Introduction}
How to structure:\\
\begin{itemize}
	\item Intoduction and review on past work on indirect reciprocity and the evolution of cooperation
	\item Methods: 
	\begin{itemize}
		\item My model of reciprocity, it's variables etc.
		\item System design (Prolog + environ)
		\item Interesting programming techniques
		\item Software engineering techniques, tools and processes
	\end{itemize}
	\item Results
	\begin{itemize}
		\item Analysis and evaluation of the results produced by my model
	\end{itemize}
	\item Discussion and Conclusion
	\begin{itemize}
		\item Impacts in real life and intelligent agents
		\item Limitations of the system and how this could be improved
		\item Issues in the development process and how these could have been avoided/accounted for
	\end{itemize}
	\item Professional issues
\end{itemize}

\chapter{Review of Literature}
Evol Coop ~\cite{evolution_of_cooperation}
Milestone: The report should describe the theory behind indirect reciprocity and it’s strategies in relation to game-theory\\
Milestone: The report should describe the link between indirect reciprocity theory and real life biological and intelligent agent interactions\\


\chapter{Contents and Knowledge}
Methods
Milestone: The report should describe the onlookers and gossip aspects of indirect reciprocity\\
Milestone: The report should contain a design of the web application and environment, Prolog service and agents, and the connection between them\\
Milestone: The report should provide an analysis and evaluation of strategies in indirect reciprocity tournaments\\
Milestone: The report should contain a discussion of the software engineering techniques, tools and processes used and issues encountered\\
Milestone: The report should contain any interesting programming techniques employed to develop the final prototype\\
Talk about design choice between kostas' idea and mine.


\chapter{Discussion and conclusions}
Discussion

\addcontentsline{toc}{chapter}{Bibliography}
\bibliography{../refs.bib}{}
\bibliographystyle{plain}

\addcontentsline{toc}{chapter}{Professional Issues}
\chapter*{Professional Issues}
Professional Issues

\end{document}
